\subsection{Resistor}

\subsubsection{Definition}

A resistor is a passive two-terminal electronic component that implements electrical resistance as a circuit element \cite{wikipedia_resistor}. It obeys Ohm's Law ($V = IR$), meaning the voltage drop across it is directly proportional to the current flowing through it. Resistance is measured in ohms ($\Omega$), and resistors values are most commonly indicated by a system of color-coded bands printed on their body.

\img{lab-01/assets/oscilloscope.jpeg}{Resistor and color-band coding chart}{registor}

\subsubsection{Specifications}
The primary specifications of a resistor are its resistance value (ranging from fractions
of an ohm up to several gigaohms), its tolerance (the permissible deviation from the
nominal value, commonly $\pm1\%$, $\pm5\%$, or $\pm10\%$), and its power rating (the
maximum heat it can safely dissipate, typically $\frac{1}{8}$\,W to $2$\,W for through-hole
types and up to several hundred watts for industrial power resistors) \cite{electronics_notes:resistor}.
The temperature coefficient of resistance (TCR), expressed in parts per million per degree
Celsius (ppm/$^\circ$C), quantifies how much the resistance drifts with temperature;
carbon-film resistors exhibit TCR values of $\pm100$ to $\pm500$\,ppm/$^\circ$C, while
metal-film types achieve as low as $\pm25$\,ppm/$^\circ$C, and precision metal-foil
resistors can reach $\pm0.2$\,ppm/$^\circ$C \cite{electronics_notes:resistor}.
Resistors are available in through-hole (axial lead) and surface-mount (SMD) packages;
SMD types are labeled using a three- or four-digit numeric code rather than color bands.